\subsection{Current control validation}

\subsubsection{Time-Domain Validation}
The time-domain validation of the current controller was performed by applying a step reference equal to the equilibrium current.
\begin{figure}[H]
	\centering
	\includegraphics[width=0.9\textwidth]{../img/current_control/time/new_FB_Current_Control_TimeValidation_CurrentTracking.png}
	\caption{Time-domain validation of the current controller.}
	\label{fig:current_time_validation}
\end{figure}


The response exhibits a fast rise time but also periodic oscillations in steady state,
attributed to mechanical disturbances present in the system.
This behavior did not compromise closed-loop performance: during position controller validation,
the current loop successfully tracked the time-varying reference generated by the outer loop, as shown in \figref{fig:fb_position_time_current}.

A potential improvement would be to reduce the current loop bandwidth and redesign the position controller accordingly,
which may reduce sensitivity to the mechanical disturbance at the cost of slower current tracking.

\subsubsection{Frequency-Domain Validation}
The frequency-domain validation was performed by applying sinusoidal references around the
equilibrium current. The tested frequencies were $100$, $500$ and $600~\mathrm{rad/s}$.
Due to the high current-loop bandwidth and the limited sampling rate, only the $100~\mathrm{rad/s}$
test provided reliable data for quantitative analysis.
\begin{figure}[H]
	\centering
	\includegraphics[width=0.9\textwidth]{../img/current_control/frequency/new_FB_Current_Control_Freq_Validation_100rads.png}
	\caption{Frequency-domain validation at $100\,\mathrm{rad/s}$.}
	\label{fig:current_freq_validation_100}
\end{figure}

The magnitude and phase variation extracted from the measured data are reported in
Table~\ref{tab:current_freq_validation}.
\begin{table}[H]
    \centering
    \renewcommand{\arraystretch}{1.6}
    \resizebox{\textwidth}{!}{%
    \begin{tabular}{|c|c|c|c|c|}
        \hline
        \textbf{\begin{tabular}[c]{@{}c@{}}Sine Wave Frequency\\ $[\mathrm{rad/s}]$\end{tabular}} &
        \textbf{\begin{tabular}[c]{@{}c@{}}Sine Wave Frequency\\ $[\mathrm{Hz}]$\end{tabular}} &
        \textbf{\begin{tabular}[c]{@{}c@{}}Cutoff Frequency\\ $[\mathrm{Hz}]$\end{tabular}} &
        \textbf{\begin{tabular}[c]{@{}c@{}}Magnitude\\ $[\mathrm{dB}]$\end{tabular}} &
        \textbf{\begin{tabular}[c]{@{}c@{}}Phase\\ $[^\circ]$\end{tabular}} \\
        \hline
        $100$ & $15.923$ & $25$ & $-0.11$ & $-11.46$ \\
        \hline
    \end{tabular}%
    }
    \caption{Frequency-domain validation results at 100 rad/s.}
    \label{tab:current_freq_validation}
\end{table}
The result is consistent with the Bode diagram of the current loop shown in \figref{fig:bode_hi}.
The sampling limitation also suggests that a lower current-loop bandwidth could simplify
experimental frequency-domain validation.
